\subsection{Übergang zu modernen Atommodellen}
\label{subsec:uebergang_moderne_atommodelle}

Die quantenmechanische Beschreibung des Elektrons\index{Elektron} verändert das
Bild des Atoms grundlegend. Im klassischen Denken erscheint das Elektron als
kleines geladenes Teilchen, das sich auf einer bestimmten Bahn um den Atomkern
\index{Atomkern} bewegt. Dieses Bild ist anschaulich, aber nur begrenzt richtig.

Bereits der Wellencharakter des Elektrons\index{Wellencharakter!des Elektrons}
zeigt, dass eine reine Teilchenbahn nicht ausreicht. Die de-Broglie-Hypothese
\index{de-Broglie-Hypothese} ordnet dem Elektron eine Wellenlänge zu. Der
Doppelspaltversuch\index{Doppelspaltversuch} zeigt, dass Elektronen
Interferenzerscheinungen\index{Interferenz} zeigen können. Damit wird deutlich:
Das Elektron besitzt keine klassische Bahn im Sinne eines kleinen Planeten um
die Sonne.

Die Schrödingergleichung\index{Schrödingergleichung} liefert die mathematische
Grundlage für dieses neue Bild. Sie beschreibt den Zustand des Elektrons durch
eine Wellenfunktion\index{Wellenfunktion}. Aus dieser Wellenfunktion ergibt
sich keine feste Bahn, sondern eine Wahrscheinlichkeitsverteilung
\index{Wahrscheinlichkeitsverteilung} für den Aufenthaltsort des Elektrons.

Damit entsteht die Orbitalvorstellung\index{Orbitalvorstellung}. Ein Orbital
\index{Orbital} ist kein Kreis und keine Bahn, sondern ein räumlicher Bereich,
in dem ein Elektron mit hoher Wahrscheinlichkeit gefunden werden kann. Die Form
und Energie dieser Orbitale ergeben sich aus der quantenmechanischen Lösung des
Atoms.

Das bohrsche Atommodell\index{bohrsches Atommodell} war ein wichtiger
Zwischenschritt. Es führte quantisierte Energieniveaus\index{Energieniveau}
ein und konnte das Spektrum des Wasserstoffatoms\index{Wasserstoffatom}
teilweise erklären. Seine Vorstellung fester Kreisbahnen wurde jedoch durch die
Quantenmechanik ersetzt.

Im modernen Atommodell\index{Atommodell!modernes} wird die Elektronenhülle
\index{Elektronenhülle} durch Quantenzahlen\index{Quantenzahl}, Orbitale und
Besetzungsregeln beschrieben. Die Hauptquantenzahl \(n\)\index{Hauptquantenzahl}
ordnet die Zustände grob nach Schalen\index{Schale!Atom}. Die Nebenquantenzahl
\(l\)\index{Nebenquantenzahl} beschreibt die Form des Orbitals. Die magnetische
Quantenzahl \(m_l\)\index{magnetische Quantenzahl} beschreibt seine räumliche
Orientierung. Die Spinquantenzahl \(m_s\)\index{Spinquantenzahl} beschreibt die
Spinprojektion\index{Spinprojektion} des Elektrons.

Zusammen mit dem Pauli-Prinzip\index{Pauli-Prinzip} ergibt sich daraus die
geordnete Struktur der Elektronenhülle. Ein Orbital kann höchstens zwei
Elektronen aufnehmen. Diese beiden Elektronen unterscheiden sich in ihrer
Spinquantenzahl. Dadurch können Elektronen nicht beliebig denselben Zustand
besetzen, sondern müssen sich auf verschiedene Zustände verteilen.

\begin{PhysicsBox}[Modernes Atommodell]
	\label{box:modernes_atommodell}
	Im modernen Atommodell bewegen sich Elektronen nicht auf festen Bahnen um den
	Atomkern. Sie werden durch Wellenfunktionen, Orbitale, Quantenzahlen und
	Besetzungsregeln beschrieben. Die Elektronenhülle entsteht aus der geordneten
	Besetzung dieser quantenmechanischen Zustände.
\end{PhysicsBox}

Diese Sichtweise erklärt die Schalenstruktur\index{Schalenstruktur} der Atome
und bildet die Grundlage für das Periodensystem\index{Periodensystem}. Elemente
mit ähnlicher äußerer Elektronenbesetzung besitzen ähnliche chemische
Eigenschaften\index{chemische Eigenschaft}. Damit wird verständlich, warum sich
die Eigenschaften der Elemente periodisch wiederholen.

Der Übergang vom bohrschen Atommodell zum modernen Orbitalmodell ist daher mehr
als eine kleine Korrektur. Er ist ein grundlegender Wechsel der Vorstellung:
Das Elektron ist nicht mehr ein Teilchen auf einer Bahn, sondern ein
quantenmechanischer Zustand mit messbaren Eigenschaften und
Wahrscheinlichkeitsverteilungen.

Damit schließt sich der gedankliche Weg dieses Kapitels. Aus dem klassischen
Elektron wurde ein quantenmechanisches Objekt\index{quantenmechanisches Objekt}.
Dieses Bild ist die Grundlage für das moderne Verständnis von Atomen,
chemischen Bindungen\index{chemische Bindung}, Festkörpern\index{Festkörper}
und elektronischen Anwendungen\index{Anwendung!elektronische}.